\chapter{L'Eterno Ritorno della Natura}

\epigraph{``L'anima \`e immortale ed \`e nata molte volte, e ha visto tutte
le cose sia qui che nell'aldil\`a.''}{\textcite{plato1997complete}, \emph{Repubblica} 614d}

\section{Il Mito di Er}

La \emph{Repubblica} termina con un mito. Er, un soldato che si credeva
morto, rinasce e racconta ci\`o che ha visto nell'aldil\`a. Descrive il
giudizio delle anime e, soprattutto, la scelta delle vite --- il momento
in cui ogni anima, prima di rinascere, sceglie il modello della vita che
vivr\`a. Alcuni scelgono saggiamente, altri scioccamente, ma tutti
scelgono secondo il carattere che hanno sviluppato nelle loro vite
precedenti.

Il mito non \`e un resoconto letterale dell'aldil\`a ma un insegnamento
filosofico sulla struttura della scelta. L'anima sceglie la sua vita, ma
la scelta non \`e arbitraria --- \`e condizionata da ci\`o che l'anima
\`e diventata. L'anima del tiranno sceglie una vita da tiranno; l'anima
del filosofo sceglie una vita da filosofo. Il mito \`e il modo di Platone
di dire che il carattere che si sviluppa in questa vita \`e il destino
che si porta nella successiva.

\section{Il Campo di Tutte le Azioni Possibili}

Riformuliamo il Mito di Er come una verit\`a strutturale su \praxis\ e
\physis. Lo spazio di embedding del motore \physis\ --- lo spazio
vettoriale ad alta dimensione in cui ogni azione, ogni intenzione, ogni
risultato \`e rappresentato --- \`e il campo di tutte le azioni possibili.
Ogni azione che il praticante potrebbe compiere occupa un punto in questo
spazio. Ogni azione che ha compiuto ne occupa un altro. La distanza tra
di essi \`e la misura dell'attualit\`a del praticante rispetto al suo
potenziale.

La scelta delle vite nel mito corrisponde alla scelta degli obiettivi
nella piattaforma \praxis. Ogni praticante, in ogni momento, sta davanti
al campo delle azioni possibili e ne sceglie una da attualizzare. Traccia
una distinzione --- questo, non quello --- e la registra come un nodo nel
grafo di coerenza. Il nodo \`e la sua traccia nel campo, il suo segno
nello spazio delle possibilit\`a.

\section{Il Giudizio di Coerenza}

Nel mito, le anime sono giudicate per le loro azioni e ricevono ci\`o che
meritano. Nel nostro quadro, il giudizio non \`e amministrato da un
giudice divino ma dalla struttura della natura stessa. Il punteggio di
coerenza \`e il giudizio: +1 per l'azione che ha prodotto un effetto
reale, 0 per l'azione che non ha lasciato traccia, -1 per l'azione che
ha violato la propria intenzione.

Questo giudizio non \`e una punizione o una ricompensa. \`E una
\emph{registrazione} --- la registrazione di ci\`o che l'azione del
praticante \`e realmente stata, distinto da ci\`o che il praticante
intendeva o credeva. Il grafo di coerenza \`e il libro della vita, ma
\`e scritto non da uno scriba esterno ma dalle azioni stesse del
praticante, nella stessa struttura del mondo.

\section{La Scelta della Prossima Azione}

Il mito descrive le anime che scelgono le loro prossime vite. Il
praticante, dopo aver rivisto il proprio grafo di coerenza, sceglie la
sua prossima azione. La scelta \`e informata dalla registrazione: il
praticante pu\`o vedere quali tipi di azioni producono coerenza per lui
e quali no. Pu\`o vedere i pattern delle proprie scelte passate.

Ma la scelta non \`e determinata dal passato. La funzione di sogno del
motore \physis\ genera continuamente variazioni nuove --- proposte per
azioni che non sono mai state provate. Il praticante non \`e condannato
a ripetere i propri pattern passati; pu\`o sempre scegliere il nuovo,
l'inprovato, l'inaspettato. Questa \`e la libert\`a che il mito preserva:
l'anima sceglie, e l'anima pu\`o scegliere diversamente.

\section{L'Immortalit\`a Attraverso la Coerenza}

Il mito conclude con l'immortalit\`a dell'anima. Offriamo un diverso tipo
di immortalit\`a: la persistenza dell'azione coerente attraverso il grafo
olarchico.

Ogni azione che ottiene +1 stabilisce un nodo nel grafo di coerenza che
pu\`o servire come modello per azioni future --- non solo per il
praticante che ha agito ma per qualsiasi praticante che acceda al grafo.
L'azione del praticante che si \`e allenato bene, che ha amato bene, che
ha costruito bene --- questa azione diventa un nodo a cui altri possono
connettersi, da cui imparare, su cui costruire.

La wiki comunitaria della piattaforma \praxis\ serve a questa funzione:
la registrazione di ci\`o che ha funzionato \`e condivisa tra i
praticanti, in modo che la coerenza di uno possa informare l'azione di un
altro. Il nodo non muore con il praticante; persiste finch\'e persiste il
grafo, finch\'e ci sono praticanti che possono connettersi ad esso.
Questa \`e l'immortalit\`a non come durata eterna ma come \emph{rilevanza
eterna} --- l'azione che continua a informare l'azione molto tempo dopo
che l'attore \`e scomparso.

\subsection{Kleos: Gloria Immortale tra i Mortali}

\`E pi\`u antico di quanto suoni. \`E il pi\`u arcaico dei desideri umani, quello
che si muove sotto l'\emph{Iliade}: \emph{kleos} --- gloria, ``ci\`o che si ode''
--- la fama immortale che l'eroe conquista quando la sua impresa \`e ricordata e
ri-cantata \parencite[IX]{homer2011iliad}. Achille sceglie una vita breve con
\emph{kleos} su una lunga senza, perch\'e l'impresa ricordata sopravvive al corpo
che l'ha compiuta. Il guerriero omerico aveva una sola tecnologia
dell'immortalit\`a: il canto dell'aedo, che portava il nome avanti solo finch\'e
v'erano cantori a portarlo. L'UberWiki \`e quel canto reso durevole. L'atto
coerente, depositato come nodo, \`e cantato non da un aedo che pu\`o dimenticare
ma da una struttura che non dimentica --- e ri-cantato ogni volta che un altro
praticante vi si connette, vi costruisce sopra, lo mette di nuovo in atto.

C'\`e una logica pi\`u profonda qui, ed \`e la pi\`u antica intuizione
aristocratica di tutte: \emph{agathoi toi agathoi} --- il buono al buono,
l'eccellenza attratta dall'eccellenza. In un bene comune governato dalla
selezione, l'istanza pi\`u coerente di un tipo-d'azione accumula il maggior numero
di connessioni; diventa il nodo che gli altri cercano per primo; e nel diventarlo,
\emph{la migliore istanza diventa la definizione operante della cosa}. L'atto pi\`u
fine di insegnare, di costruire, di curare che il grafo contiene non \`e pi\`u un
insegnamento tra molti --- diventa ci\`o che insegnare, qui, \emph{significa}: il
riferimento contro cui il prossimo \`e letto, la Forma che il migliore
particolare si \`e guadagnato il diritto di essere. \`E la strana immortalit\`a che
il praticante pu\`o conquistare. Non che il suo nome sia inciso, ma che la sua
impresa, essendo stata la pi\`u coerente, \emph{diventi lo standard} --- sicch\'e
chiunque agisca bene dopo di lui sta, in parte, mettendo in atto lui. L'\emph{agathos}
non consegue meramente l'eccellenza; l'atto eccellente, depositato, attrae a s\'e
il prossimo \emph{agathos}, e l'eccellenza si propaga all'eccellenza attraverso il
bene comune. Il miglior processo diventa la definizione del fine a cui mirava.

Cos\`i il mortale conquista, attraverso l'UberWiki, esattamente ci\`o che l'eroe
omerico voleva e poteva assicurarsi solo per fortuna degli aedi: che la sua azione
pi\`u fine non passi. \`E assunta nella struttura di ci\`o che altri faranno, un
nodo immortale in un grafo che continua a ri-cantarlo. Questa \`e gloria tra i
mortali --- non la fama vuota del regime timocratico (Libro VIII), che cercava di
essere \emph{visto}, ma il vero \emph{kleos} dell'\emph{agathos}, che cercava di
agire bene e, agendo bene, divenne la forma dell'atto stesso. Il vir spende il suo
\textsc{ayu} in atto coerente; l'atto diventa la definizione; la definizione
sopravvive al vir. \`E tutto qui.

\begin{figure}[ht]
\centering
\begin{tikzpicture}[
  mynode/.style={draw, circle, minimum size=6mm, inner sep=1pt, font=\tiny},
  >={Stealth[length=1.5mm]}]
  \foreach \ang in {0,60,120,180,240,300}
    \node[mynode, fill=cohgreen!30, draw=cohgreen] (g\ang) at (\ang:10mm) {};
  \node[mynode, fill=cohgreen!50, draw=cohgreen, minimum size=7mm] (gc) at (0,0) {};

  \foreach \ang in {0,90,180,270}
    \node[mynode, fill=cohgray!30, draw=cohgray] (y\ang) at (\ang:22mm) {};

  \node[mynode, fill=cohred!30, draw=cohred] (r1) at (135:28mm) {};
  \node[mynode, fill=cohred!30, draw=cohred] (r2) at (160:30mm) {};

  \draw[cohgreen!40] (gc) -- (g0) (gc) -- (g60) (gc) -- (g120)
    (gc) -- (g180) (gc) -- (g240) (gc) -- (g300);
  \draw[cohgreen!30] (g0) -- (g60) (g60) -- (g120) (g120) -- (g180)
    (g180) -- (g240) (g240) -- (g300) (g300) -- (g0);
  \draw[cohgray!40] (g0) -- (y0) (g60) -- (y90) (g180) -- (y270) (g240) -- (y180);

  \draw[->, dashed, cohred!60] (r1) to[bend left] ++(10:10mm);
  \draw[->, dashed, cohred!60] (r2) to[bend right] ++(190:8mm);

  \node[font=\tiny, cohgreen, above] at (gc.north) {Coerente (+1)};
  \node[font=\tiny, cohgray, above] at (y0.north) {Inerte (0)};
  \node[font=\tiny, cohred, above] at (r1.north) {Fallimento (--1)};
  \node[font=\tiny, align=center, below=2mm] at (current bounding box.south)
    {Frecce tratteggiate: funzione sogno che propone variazioni sui nodi falliti.};
\end{tikzpicture}
\caption{Il grafo olarchico delle azioni. I nodi verdi sono azioni coerenti che
persistono; i nodi grigi sono inerti; i nodi rossi sono fallimenti riprodotti dalla
funzione sogno.\label{fig:olarchia}}
\end{figure}

\section{La Visibilit\`a del S\'e-Impersonale}

Prima dell'argomento conclusivo, una mossa preparatoria. Il grafo olarchico fa
qualcosa a \emph{das Man} --- l'anonimo ``si'' di Heidegger, lo standard pubblico
impersonale di ``ci\`o che si fa'' che governa inautenticamente proprio perch\'e
nessuno lo firma ed \`e dappertutto, l'aria che si respira
\parencite[§27]{heidegger1962being}. Il grafo lo \emph{reifica}: proiettato in uno
spazio vettoriale condiviso, il s\'e-impersonale appare come una struttura di nodi
pesati per coerenza, e ci\`o che era l'orizzonte invisibile dell'impegno diventa
una cosa che si pu\`o guardare. La reificazione, qui, non \`e alienazione ma il suo
opposto --- poich\'e ci\`o che pu\`o essere visto non pu\`o pi\`u governare non
visto.

Il grafo rende visibile il s\'e-impersonale dandogli le tre cose dietro cui si
nascondeva. Un \emph{tempo}: il fallimento ricorrente appare come un ritmo, non
una serie di accidenti, e un ritmo una volta visto pu\`o essere spezzato. Una
\emph{causa}: ogni nodo porta un bordo verso il soggetto che lo ha registrato,
sicch\'e l'atto ``che nessuno ha deciso ma tutti seguono'' \`e smascherato come
catena di responsabilit\`a abbandonata. E un \emph{luogo}: la distanza tra il
proprio vettore e il centro di gravit\`a statistico di ogni azione \`e una distanza
misurabile dall'inautenticit\`a --- ci\`o che era dappertutto e in nessun luogo
diventa localizzabile. Vedere cos\`i il s\'e-impersonale \`e la prima met\`a del
sottrarsi ad esso. La seconda met\`a --- renderlo non solo visibile ma
\emph{responsabile} --- \`e il lavoro della repubblica a cui ora ci volgiamo.

\section{La Coerenza dell'Intero}

La \emph{Repubblica} termina con un'immagine di ordine cosmico: il fuso
della Necessit\`a che gira le sfere del cosmo, le Sirene che cantano
ciascuna una singola nota, le tre Parche che tessono il destino di ogni
anima nel tessuto dell'intero. Le scelte individuali contano, ma contano
all'interno di un cosmo che ha il proprio ordine, la propria musica, la
propria coerenza.

Il grafo olarchico del motore \physis\ \`e l'immagine di questo ordine
cosmico al livello dell'azione. Ogni nodo \`e connesso ad altri da bordi
di influenza, causalit\`a e risonanza. L'azione del praticante non \`e
mai isolata; si increspa attraverso il grafo, influenzando nodi che il
praticante potrebbe non conoscere mai. La coerenza dell'intero non \`e
la somma delle coerenze individuali ma la struttura della loro
interconnessione --- la forma del grafo stesso
\parencite{koestler1967ghost}.

\section{La Repubblica dal Basso}

Abbiamo seguito la \emph{Repubblica} di Platone per dieci libri, e alla fine
dobbiamo confessare l'unico punto in cui non lo seguiamo affatto --- ed \`e il
punto pi\`u importante. La repubblica di Platone \`e costruita dall'alto verso il
basso. Il filosofo ascende solo, vede il Bene, e \emph{discende a governare}:
l'autorit\`a fluisce verso il basso da colui che ha visto ai molti che non hanno
visto. La citt\`a \`e giusta quando i saggi comandano e il resto obbedisce. \`E
l'architettura che ogni lettore dall'antichit\`a ha trovato o esaltante o
mostruosa, e poggia su un'unica assunzione portante: che il Bene debba essere
\emph{visto}, da pochi, e poi \emph{imposto} ai molti.

La nostra repubblica inverte questo. Non c'\`e praticante-re, perch\'e non c'\`e
nessuno che veda il Bene e governi col vederlo. Il Bene --- che per noi \`e
\physis, coerenza, il verdetto della natura sull'azione --- non \`e visto
dall'alto e tramandato. \`E \emph{scoperto frammento per frammento, dal basso},
da ciascun praticante nella grana della propria azione, e \emph{depositato} in un
deposito comune. Quel deposito comune \`e l'UberWiki: l'unico grafo olarchico
condiviso in cui le azioni coerenti di ogni praticante sono scritte come nodi che
ogni altro praticante pu\`o trovare, collegare e su cui costruire. L'autorit\`a
non discende; si \emph{accumula}.

L'UberWiki \`e, nel termine di Bernard Stiegler, un sistema di \emph{ritenzione
terziaria}: memoria esteriorizzata in un supporto durevole fuori dall'individuo,
il terzo tipo di ritenzione oltre il presente vissuto (primaria) e il ricordo
personale (secondaria) \parencite{stiegler1998technics}. L'intuizione di
Stiegler, seguendo il mito di Epimeteo, \`e che l'umano \`e l'essere costituito
dalle proprie protesi --- che non abbiamo natura se non i supporti tecnici
attraverso cui ereditiamo un passato che non abbiamo vissuto e trasmettiamo un
futuro che non vedremo. \`E a doppio taglio, e Stiegler lo sapeva: la stessa
memoria esteriorizzata che rende possibile la cultura \`e, nelle mani
dell'economia dell'attenzione, lo strumento della sua cattura --- ritenzione
industrializzata, passato pre-confezionato, desiderio cortocircuitato. La domanda
decisiva non \`e dunque se esteriorizzare la memoria --- non possiamo non farlo
--- ma \emph{sotto quale governatore} il sistema di ritenzione operi. Il feed \`e
ritenzione terziaria governata dall'engagement; l'UberWiki \`e ritenzione
terziaria governata dalla coerenza. Entrambi sono protesi della memoria; uno
cattura il praticante, l'altro lo capacita. \`E il senso preciso in cui la
repubblica dal basso non \`e una fuga dalla tecnologia ma la sua redenzione: non
abolisce la memoria esteriorizzata che ci rende umani, la ri-subordina al
verdetto della natura anzich\'e a quello del mercato.

Non \`e una metafora imposta al sistema a posteriori. \`E l'architettura
letterale del motore in esecuzione. Ogni praticante agisce nel proprio stato
isolato --- i propri obiettivi, il proprio ciclo PDCA, il proprio ciclo del
sogno, il proprio nucleo di coerenza, sigillati da quelli di ogni altro
praticante. Il lavoro privato del divenire \`e privato. Ma il grafo olarchico \`e
\emph{globale}: un'unica struttura di conoscenza condivisa, alimentata da un solo
condotto di ingestione, comune a tutti. Ci\`o che \`e privato \`e lo sforzo;
ci\`o che \`e condiviso \`e il risultato. Il praticante fatica solo e contribuisce
a un bene comune. L'UberWiki \`e la citt\`a costruita non nel discorso, come
Socrate costru\`i la sua, ma nella \emph{coerenza depositata} --- il sedimento di
ogni azione che ha realmente funzionato, lasciata dove il prossimo praticante
pu\`o starci sopra. E questo bene comune non \`e confinato a una sola macchina: il
grafo di coerenza \`e rispecchiato in un deposito durevole e interrogabile fuori
da ogni singolo dispositivo, sicch\'e la citt\`a della coerenza depositata
sopravvive alla sessione che l'ha scritta e pu\`o essere condivisa tra i
praticanti che la tengono in comune --- alla repubblica \`e data non solo una
costituzione ma un territorio.

In una tale repubblica, il governo non \`e comando ma \emph{selezione}. Nessun
nodo governa un altro per decreto. Un nodo guadagna autorit\`a come un sentiero
la guadagna in una foresta o una citazione in una letteratura: per essere usato,
collegato, percorso di nuovo. L'azione che ha prodotto coerenza reale diventa un
nodo a cui altri si collegano; accumula bordi; diventa portante nel grafo non
perch\'e qualcuno l'ha ordinato ma perch\'e ha continuato a dimostrarsi. Il nodo
incoerente deriva inerte, non collegato, e infine \`e rigiocato e sovrascritto
dalla funzione del sogno. \`E la stessa logica con cui la natura governa senza un
governatore: ritenzione selettiva di ci\`o che funziona, cancellazione silenziosa
di ci\`o che non funziona. L'UberWiki \`e governato da nessuno e ordinato da
tutti --- una repubblica il cui unico legislatore \`e il registro accumulato di
ci\`o che ha realmente dato frutto.

Due conseguenze seguono, e sono la soluzione di due debiti contratti prima in
quest'opera. Primo, il praticante-re del Libro V non \`e abolito ma
\emph{universalizzato}. Platone aveva un re e molti sudditi; noi non abbiamo
sudditi, perch\'e ogni praticante \`e re della propria azione e governato solo
dalla coerenza con la forma che egli stesso ha dichiarato. La regalit\`a del
Libro V non \`e un rango a cui pochi ascendono ma una condizione che ogni
praticante gi\`a occupa nel momento in cui assume la propria entelechia. La
repubblica non \`e una gerarchia di uno sui molti ma una federazione di sovrani,
ciascuno assoluto nel proprio dominio, ciascuno che contribuisce i propri
risultati verificati al grafo comune.

Secondo, l'UberWiki rende \emph{das Man} non solo visibile ma \emph{responsabile}.
Vedere il s\'e-impersonale, come ha mostrato la sezione precedente, \`e sottrarsi
inconsapevolmente al suo dominio; ma la sola visibilit\`a lascia il grafo
condiviso capace di divenire la camera pi\`u profonda della caverna ---
l'aggregato statistico di ``ci\`o che si fa,'' un algoritmo che governa
precisamente perch\'e nessuno l'ha scritto. La stessa struttura condivisa \`e o la
tirannia dell'anonimo o il senato dei responsabili, e tutta la differenza sta
nella \emph{paternit\`a}. Nel
feed dell'economia dell'attenzione, l'aggregato non ha autori: nessun nodo possiede
la propria origine, nessun contributo porta un nome, il ``si'' parla e nessuno
\`e responsabile. Nell'UberWiki, ogni nodo ha un soggetto registrato e una catena
causale tracciabile. L'aggregato \`e composto di contributi \emph{firmati,
posseduti, responsabili}. \`E ci\`o che converte \emph{das Man} in una
repubblica: non l'abolizione dell'orizzonte condiviso --- \`e impossibile, siamo
sempre gi\`a in un mondo comune --- ma la sua \emph{ri-autorizzazione}, il
ripristino di un nome e di una causa a ogni nodo del bene comune. La repubblica
dal basso \`e il s\'e-impersonale reso responsabile verso se stesso.

L'ordine cosmico con cui Platone chiude --- le sfere, le Sirene, le Parche --- era
una gerarchia fissata dall'alto, ogni anima intessuta in un posto assegnato dalla
Necessit\`a. Il nostro ordine cosmico \`e tessuto dal basso: un grafo senza centro
e senza sovrano, la cui coerenza \`e il precipitato vivente di ogni praticante che
abbia mai agito bene e lasciato la traccia dove altri potessero trovarla. \`E una
repubblica non perch\'e qualcuno l'ha progettata giusta, ma perch\'e la natura la
governa, e la natura governa dal basso.

\section{Il Ritorno all'Inizio}

Finiamo dove abbiamo iniziato: con l'affermazione che \praxis\ \`e lo
strumento di \physis, e che il governo dell'azione da parte della natura
\`e la condizione dell'azione retta. I dieci libri hanno tracciato
l'argomento attraverso l'omologia di natura e azione, le discipline del
praticante, la natura della coerenza come giustizia, l'ideale del
praticante-re, la linea divisa dell'azione, la fuga dalla caverna della
Trappola dello Specchio, il declino di \praxis\ nelle sue forme corrotte
e la risubordinazione della tecnologia alla natura.

L'argomento \`e completo, ma il lavoro no. Il lavoro \`e la vita del
praticante --- la disciplina quotidiana del check-in, della testimonianza,
della posta in gioco, della revisione. Il lavoro \`e la costruzione e la
manutenzione del motore \physis\ e della piattaforma \praxis\ --- la
prova per costruzione che la filosofia pu\`o diventare architettura. Il
lavoro \`e il ritorno alla caverna, ancora e ancora, per testimoniare per
coloro che sono ancora incatenati.

La \emph{Repubblica} termina con una promessa: ``Se saremo persuasi da me,
manterremo salda la via verso l'alto, e praticheremo la giustizia con
saggezza, in modo da essere amici di noi stessi e degli d\`ei.''
Offriamo una promessa corrispondente: se saremo persuasi dall'argomento
di questa tesi, manterremo salda la via della coerenza. Praticheremo
l'azione con sintonia alla natura. Diventeremo amici di noi stessi e
dell'intero, di cui siamo gli strumenti pi\`u articolati.
